\documentclass[11pt,a4paper]{article}

\usepackage[T1]{fontenc}
\usepackage[utf8]{inputenc}
\usepackage{lmodern}
\usepackage{amssymb,amsmath,amsthm,mathtools}
\usepackage[margin=2.5cm]{geometry}
\usepackage{parskip}
\usepackage[hidelinks]{hyperref}
\usepackage{microtype}

% Theorem environments (arXiv / journal style)
\theoremstyle{plain}
\newtheorem{theorem}{Theorem}[section]
\newtheorem{proposition}[theorem]{Proposition}
\newtheorem{lemma}[theorem]{Lemma}
\newtheorem{corollary}[theorem]{Corollary}

\theoremstyle{definition}
\newtheorem{definition}[theorem]{Definition}
\newtheorem{example}[theorem]{Example}
\newtheorem{remark}[theorem]{Remark}

% Notation
\newcommand{\Type}{\mathbf{Type}}
\newcommand{\bott}{\mathrm{bot}}
\newcommand{\cmin}{\mathrm{cmin}}
\newcommand{\cmax}{\mathrm{cmax}}
\newcommand{\tauCan}{\tau}
\newcommand{\Prem}{\mathrm{Prem}}
\newcommand{\Sound}{\mathrm{Sound}}
\newcommand{\Deriv}{\mathrm{Deriv}}

\title{Trust Systems:\\
  Canonical Trust, Certificates, and Consolidation}
\author{Extracted from a machine-checked development}
\date{}

\begin{document}
\maketitle

\begin{abstract}
We develop the elementary theory of \emph{trust systems}: dependency
structures on claims, graded in a linearly ordered poset of check
strengths. Each system carries a \emph{canonical trust} assignment~$\tau$,
characterised both as the greatest sound assignment of grades to claims
and as the supremum of grades of derivations. Morphisms of trust systems
transport canonical trust laxly. \emph{Consolidations}---shape maps
equipped with position pullbacks that cache grades---preserve~$\tau$ on
the nose. A circular-citation counterexample shows that well-foundedness
of the premise relation is load-bearing; a shape-only consolidation
counterexample shows that dropping checking steps degrades trust.
All stated results have been machine-checked.
\end{abstract}

\tableofcontents
\bigskip

%======================================================================
\section{Introduction}
%======================================================================

A recurring situation in formal and informal mathematical work is that a
body of claims is justified by checking steps, each step resting on
premises and receiving a grade recording how thoroughly it was checked.
One wants a coherent global notion of how much each claim can be trusted,
computed from the local data, and one wants to know how that global grade
moves when claims are restated, imported, or consolidated across layers.

Trust systems make this precise. A trust system is a type of claims, a
type of checking steps, and a grading of steps in a chain of check
strengths. The \emph{canonical trust}~$\tau$ of a claim is the
widest-bottleneck value: the best derivation's weakest link. It may be
defined either by well-founded recursion on the premise relation
(fixpoint form) or as the supremum of derivation grades (max--min form);
the two agree.

\paragraph{What this note claims.}
In ordinary mathematical language: every sound assignment of grades is
bounded by~$\tau$, and~$\tau$ itself is sound, so~$\tau$ is the greatest
sound assignment. Morphisms transport trust laxly
($\phi(\tau(x)) \le \tau'(F(x))$). A consolidation with a genuine
position map preserves trust exactly
($\tau'(F(x)) = \tau(x)$). Without well-foundedness of the premise
relation, local certificates can certify circular trust that no
derivation supports; without position pullbacks, restating claims alone
can drop trust from the top grade to bottom.

%======================================================================
\section{Chains of check strengths}
%======================================================================

\begin{definition}[Chain]\label{def:chain}
A \textbf{chain} is a type~$T$ equipped with a preorder~$\le$ such that
\begin{enumerate}
  \item $\le$ is a partial order (reflexive, transitive, antisymmetric);
  \item $\le$ is total: for all $a,b \in T$, either $a \le b$ or $b \le a$;
  \item there is a bottom element $\bott \in T$ with $\bott \le a$ for every
        $a \in T$;
  \item $\le$ is decidable.
\end{enumerate}
We write $\bott$ for ``no check performed''.
\end{definition}

\begin{definition}[Meet and join]\label{def:cmin-cmax}
On a chain~$T$, define the binary meet and join by
\begin{align}
  \cmin(a,b) &\;=\;
  \begin{cases}
    a & \text{if } a \le b, \\
    b & \text{otherwise},
  \end{cases}
  \\[4pt]
  \cmax(a,b) &\;=\;
  \begin{cases}
    b & \text{if } a \le b, \\
    a & \text{otherwise}.
  \end{cases}
\end{align}
\end{definition}

\begin{proposition}\label{prop:cmin-cmax}
For all $a,b \in T$:
\begin{enumerate}
  \item $\cmin(a,b) \le a$ and $\cmin(a,b) \le b$;
  \item if $c \le a$ and $c \le b$, then $c \le \cmin(a,b)$;
  \item $a \le \cmax(a,b)$ and $b \le \cmax(a,b)$;
  \item $\cmax(a,b)$ equals~$a$ or equals~$b$.
\end{enumerate}
\end{proposition}

\begin{proof}
By lean.
\end{proof}

We extend $\cmin$ and $\cmax$ to finite lists by right folds, writing
$\mathrm{foldr}_{\cmax}(\ell, i)$ and $\mathrm{foldr}_{\cmin}(\ell, i)$
for the fold of a list~$\ell$ with initial value~$i$.

\begin{proposition}[Fold properties]\label{prop:folds}
Let $T$ be a chain.
\begin{enumerate}
  \item If $a \in \ell$, then $a \le \mathrm{foldr}_{\cmax}(\ell, i)$.
  \item $\mathrm{foldr}_{\cmax}(\ell, i)$ equals~$i$ or belongs to~$\ell$
        (attainment).
  \item If $a \in \ell$, then $\mathrm{foldr}_{\cmin}(\ell, i) \le a$;
        also $\mathrm{foldr}_{\cmin}(\ell, i) \le i$.
  \item Pointwise domination of mapped values lifts to $\cmin$-folds.
  \item Every monotone map $\phi : T \to T'$ is lax for $\cmin$-folds:
        $\phi(\mathrm{foldr}_{\cmin}(\ell, i))
         \le \mathrm{foldr}_{\cmin}(\phi(\ell), \phi(i))$.
\end{enumerate}
\end{proposition}

\begin{proof}
By lean.
\end{proof}

\begin{remark}
Attainment of finite joins (item~(2) above) uses totality of the chain:
in a general lattice a finite join need not be attained by any single
element. All check-posets used in applications are finite chains.
\end{remark}

%======================================================================
\section{Trust systems and canonical trust}
%======================================================================

\begin{definition}[Trust system]\label{def:system}
Let $C$ be a type of claims and let $T$ be a chain. A \textbf{trust
system}~$S$ over~$C$ graded in~$T$ consists of
\begin{enumerate}
  \item a type $\mathrm{Step}$ of checking steps;
  \item a conclusion map $\mathrm{concl} : \mathrm{Step} \to C$;
  \item a premise list $\mathrm{prems} : \mathrm{Step} \to \mathrm{List}(C)$;
  \item a grade map $\mathrm{grade} : \mathrm{Step} \to T$;
  \item for each claim~$x$, a finite list $\mathrm{stepsInto}(x)$ of steps
        concluding~$x$, such that
        $f \in \mathrm{stepsInto}(x)$ if and only if $\mathrm{concl}(f) = x$.
\end{enumerate}
\end{definition}

\begin{definition}[Premise relation]\label{def:prem}
The \textbf{premise relation} $\Prem_S \subseteq C \times C$ is defined by
\[
  y \;\Prem_S\; x
  \quad\Longleftrightarrow\quad
  \exists\, f \in \mathrm{stepsInto}(x)
  \text{ with } y \in \mathrm{prems}(f).
\]
Well-foundedness of $\Prem_S$ is the standing hypothesis for the
certificate lemma below.
\end{definition}

\begin{definition}[Step grade]\label{def:stepgrade}
Given an assignment $t : C \to T$, the \textbf{grade conferred by a
step}~$f$ under~$t$ is the meet of the step's own grade with the assigned
levels of its premises:
\begin{equation}
  \mathrm{stepGrade}_S(t,f)
  \;=\;
  \mathrm{foldr}_{\cmin}\bigl(
    t\bigl(\mathrm{prems}(f)\bigr),\;
    \mathrm{grade}(f)
  \bigr).
\end{equation}
\end{definition}

\begin{definition}[Canonical trust]\label{def:tau}
Assume $\Prem_S$ is well-founded. The \textbf{canonical trust}
$\tauCan_S : C \to T$ is defined by well-founded recursion on~$\Prem_S$:
\begin{equation}
  \tauCan_S(x)
  \;=\;
  \mathrm{foldr}_{\cmax}\Bigl(
    \bigl\{\,
      \mathrm{stepGrade}_S(\tauCan_S, f)
      \bigm|\,
      f \in \mathrm{stepsInto}(x)
    \,\bigr\},\;
    \bott
  \Bigr).
\end{equation}
In words: $\tauCan_S(x)$ is the best step into~$x$, valued at its step-grade
under~$\tauCan_S$ itself.
\end{definition}

\begin{definition}[Sound assignment]\label{def:sound}
An assignment $t : C \to T$ is \textbf{sound} if every claim trusted above
bottom carries a node-local certificate: for every $x$ with $t(x) \ne \bott$,
there exists $f \in \mathrm{stepsInto}(x)$ such that
\[
  t(x) \;\le\; \mathrm{stepGrade}_S(t, f).
\]
\end{definition}

\begin{theorem}[Certificate lemma, upper bound]\label{thm:sound-le-tau}
Assume $\Prem_S$ is well-founded. Every sound assignment is bounded by
canonical trust: if $t$ is sound, then
\[
  t(x) \;\le\; \tauCan_S(x)
  \qquad\text{for every claim } x.
\]
\end{theorem}

\begin{proof}
By lean.
\end{proof}

\begin{theorem}[Certificate lemma, soundness of $\tau$]\label{thm:tau-sound}
Assume $\Prem_S$ is well-founded. Canonical trust is itself sound.
\end{theorem}

\begin{proof}
By lean.
\end{proof}

\begin{corollary}[Greatest sound assignment]\label{cor:greatest}
Under well-foundedness of $\Prem_S$, the assignment $\tauCan_S$ is the
greatest sound assignment of grades to claims.
\end{corollary}

\begin{proof}
By lean.
\end{proof}

\begin{remark}
The soundness of~$\tau$ uses attainment of finite joins, hence totality of
the chain. The upper bound (Theorem~\ref{thm:sound-le-tau}) would hold more
generally.
\end{remark}

%======================================================================
\section{Derivations and the max--min semantics}
%======================================================================

\begin{definition}[Derivation]\label{def:deriv}
A \textbf{derivation} of a claim~$x$ in a trust system~$S$ is a checking
step~$f$ with $\mathrm{concl}(f) = x$, together with a derivation of each
premise of~$f$:
\begin{equation}
  \frac{
    f : \mathrm{Step}
    \qquad
    \bigl(\forall\, y \in \mathrm{prems}(f).\;
      \Deriv_S(y)\bigr)
  }{
    \Deriv_S\bigl(\mathrm{concl}(f)\bigr)
  }\,.
\end{equation}
\end{definition}

\begin{definition}[Derivation grade]\label{def:derivgrade}
The \textbf{grade of a derivation} is the meet of every step-grade in the
derivation tree (the weakest link): for a node with step~$f$ and
sub-derivations of grades $g_y$,
\begin{equation}
  \mathrm{derivGrade}\bigl(\mathrm{node}(f, \mathrm{sub})\bigr)
  \;=\;
  \mathrm{foldPrems}\bigl(
    \mathrm{prems}(f),\;
    y \mapsto \mathrm{derivGrade}(\mathrm{sub}_y),\;
    \mathrm{grade}(f)
  \bigr).
\end{equation}
\end{definition}

\begin{theorem}[Fixpoint dominates derivations]\label{thm:deriv-le-tau}
Assume $\Prem_S$ is well-founded. For every derivation~$d$ of a claim~$x$,
\[
  \mathrm{derivGrade}(d) \;\le\; \tauCan_S(x).
\]
\end{theorem}

\begin{proof}
By lean.
\end{proof}

\begin{theorem}[Attainment by a derivation]\label{thm:tau-le-deriv}
Assume $\Prem_S$ is well-founded. If $\tauCan_S(x) \ne \bott$, then there
exists a derivation~$d$ of~$x$ such that
\[
  \tauCan_S(x) \;\le\; \mathrm{derivGrade}(d).
\]
\end{theorem}

\begin{proof}
By lean.
\end{proof}

\begin{corollary}[Agreement of fixpoint and max--min forms]
\label{cor:agree}
Under well-foundedness, the fixpoint form of~$\tau$ and the
max-over-derivations of min-over-grades semantics coincide: whenever
$\tauCan_S(x) \ne \bott$, the value $\tauCan_S(x)$ is attained by some
derivation and bounds every derivation from above.
\end{corollary}

\begin{proof}
By lean.
\end{proof}

%======================================================================
\section{Morphisms and transport of trust}
%======================================================================

\begin{definition}[Morphism of trust systems]\label{def:mor}
Let $S$ be a trust system over claims~$C$ graded in a chain~$T$, and let
$S'$ be a trust system over claims~$C'$ graded in a chain~$T'$. A
\textbf{morphism} $M : S \to S'$ consists of
\begin{enumerate}
  \item a map on claims $F : C \to C'$;
  \item a map on steps $F_{\mathrm{Step}} : \mathrm{Step}_S \to
        \mathrm{Step}_{S'}$;
  \item a map on grades $\phi : T \to T'$ that is monotone,
        satisfies $\phi(\bott) = \bott$, and is lax on step grades:
        \[
          \phi\bigl(\mathrm{grade}_S(f)\bigr)
          \;\le\;
          \mathrm{grade}_{S'}\bigl(F_{\mathrm{Step}}(f)\bigr);
        \]
  \item premise preservation:
        $\mathrm{prems}_{S'}(F_{\mathrm{Step}}(f))
         = F\bigl(\mathrm{prems}_S(f)\bigr)$;
  \item step-membership preservation:
        if $f \in \mathrm{stepsInto}_S(x)$, then
        $F_{\mathrm{Step}}(f) \in \mathrm{stepsInto}_{S'}(F(x))$.
\end{enumerate}
\end{definition}

\begin{remark}
Strictness of~$\phi$ at bottom---``unchecked maps to unchecked''---is
essential. Monotonicity alone fails when a source claim has no derivation
but $\phi(\bott)$ exceeds the target's bottom.
\end{remark}

\begin{theorem}[Lax transport of trust]\label{thm:tau-transport}
Let $M : S \to S'$ be a morphism, and assume both premise relations are
well-founded. Then for every claim~$x$,
\begin{equation}
  \phi\bigl(\tauCan_S(x)\bigr)
  \;\le\;
  \tauCan_{S'}\bigl(F(x)\bigr).
\end{equation}
Importing along a morphism never claims more than the image under~$\phi$
of the source's canonical trust.
\end{theorem}

\begin{proof}
By lean.
\end{proof}

%======================================================================
\section{The circular-citation counterexample}
%======================================================================

Well-foundedness of the premise relation is load-bearing. We exhibit a
two-claim cycle in which every node-local certificate passes, yet no
derivation exists at all.

\begin{definition}[Two-level chain]\label{def:two}
Let $\mathbf{Two} = \{\mathrm{lo}, \mathrm{hi}\}$ with
$\mathrm{lo} \le \mathrm{hi}$ (and $\mathrm{lo} = \bott$).
\end{definition}

\begin{definition}[Circular system]\label{def:cyc}
Let claims be booleans $\mathbb{B} = \{\mathsf{true}, \mathsf{false}\}$.
Define a trust system $\mathrm{cyc}$ over $\mathbb{B}$ graded in
$\mathbf{Two}$ by: the unique step into a claim~$x$ has single premise
$\neg x$ and grade~$\mathrm{hi}$.
\end{definition}

\begin{proposition}[Constant-top is sound]\label{prop:cyc-sound}
The constant assignment $t(x) = \mathrm{hi}$ is sound on~$\mathrm{cyc}$.
\end{proposition}

\begin{proof}
By lean.
\end{proof}

\begin{proposition}[No derivations]\label{prop:cyc-no-deriv}
No claim in~$\mathrm{cyc}$ admits any derivation.
\end{proposition}

\begin{proof}
By lean.
\end{proof}

\begin{proposition}[Premise relation not well-founded]\label{prop:cyc-not-wf}
The premise relation of~$\mathrm{cyc}$ is not well-founded.
\end{proposition}

\begin{proof}
By lean.
\end{proof}

\begin{remark}
Circular citation is trust inflation invisible to any local check:
well-foundedness is a global hypothesis and must be checked globally. The
certificate lemma does not apply---as it must not.
\end{remark}

%======================================================================
\section{Consolidation preserves trust}
%======================================================================

Consolidation restates claims (a shape map) and, when done with provenance,
carries a position map: every checking step into a consolidated claim pulls
back to a checking step into each source claim, with corresponding
premises and a cached grade not exceeding the source grade.

\begin{definition}[Consolidation]\label{def:consolidation}
A \textbf{consolidation} $K$ from a trust system~$S$ over~$C$ to a trust
system~$S'$ over~$C'$, both graded in the same chain~$T$, consists of
\begin{enumerate}
  \item a map on claims $F : C \to C'$;
  \item a forward map on steps $F_{\mathrm{Step}} : \mathrm{Step}_S \to
        \mathrm{Step}_{S'}$ preserving premises (after applying~$F$),
        lax on grades ($S.\mathrm{grade}(f) \le
        S'.\mathrm{grade}(F_{\mathrm{Step}}(f))$), and preserving
        step-membership;
  \item a backward map on positions: for each $x \in C$ and each step
        $f' \in \mathrm{stepsInto}_{S'}(F(x))$, a step
        $\mathrm{posn}(f') \in \mathrm{stepsInto}_S(x)$ such that
        \[
          \mathrm{prems}_{S'}(f')
          \;=\;
          F\bigl(\mathrm{prems}_S(\mathrm{posn}(f'))\bigr)
        \]
        and
        \[
          S'.\mathrm{grade}(f')
          \;\le\;
          S.\mathrm{grade}\bigl(\mathrm{posn}(f')\bigr).
        \]
\end{enumerate}
The position map is indexed by the preimage~$x$: when consolidation merges
claims, every justification of the merged claim must pull back on
\emph{each} side.
\end{definition}

\begin{remark}
The definition of consolidation never mentions~$\tau$. The condition is
bisimulation-style and node-local; preservation of canonical trust is a
theorem, not a definition.
\end{remark}

\begin{proposition}[Forward half is a morphism]\label{prop:to-mor}
Every consolidation determines a morphism of trust systems at
$\phi = \mathrm{id}_T$.
\end{proposition}

\begin{proof}
By lean.
\end{proof}

\begin{lemma}[No trust inflation]\label{lem:tau-le}
Let $K$ be a consolidation from~$S$ to~$S'$, and assume both premise
relations are well-founded. Then for every claim~$x$,
\[
  \tauCan_{S'}\bigl(F(x)\bigr)
  \;\le\;
  \tauCan_S(x).
\]
\end{lemma}

\begin{proof}
By lean.
\end{proof}

\begin{theorem}[Consolidation preservation]\label{thm:tau-consolidate}
Let $K$ be a consolidation from~$S$ to~$S'$, and assume both premise
relations are well-founded. Then for every claim~$x$,
\begin{equation}
  \tauCan_{S'}\bigl(F(x)\bigr)
  \;=\;
  \tauCan_S(x).
\end{equation}
\end{theorem}

\begin{proof}
By lean.
\end{proof}

\begin{remark}
The forward inequality is lax transport at $\phi = \mathrm{id}$; the
backward inequality is Lemma~\ref{lem:tau-le}, using the position map to
pull every target justification back to the source.
\end{remark}

%======================================================================
\section{Composition of consolidations}
%======================================================================

\begin{definition}[Identity consolidation]\label{def:id-consol}
The identity consolidation on a trust system~$S$ has identity maps on
claims and steps, and identity position pullbacks, with exact (not merely
lax) grade caches.
\end{definition}

\begin{definition}[Composition of consolidations]\label{def:comp-consol}
If $K : S \to S'$ and $L : S' \to S''$ are consolidations (same chain~$T$),
their composite $K ; L : S \to S''$ has
\begin{itemize}
  \item shapes composed covariantly:
        $F_{K;L} = F_L \circ F_K$,
        $(F_{\mathrm{Step}})_{K;L} =
         (F_{\mathrm{Step}})_L \circ (F_{\mathrm{Step}})_K$;
  \item positions composed contravariantly:
        $\mathrm{posn}_{K;L}(f'') =
         \mathrm{posn}_K\bigl(\mathrm{posn}_L(f'')\bigr)$;
  \item cached grades chained by transitivity of~$\le$.
\end{itemize}
\end{definition}

\begin{theorem}[End-to-end preservation]\label{thm:comp-preserve}
If $K : S \to S'$ and $L : S' \to S''$ are consolidations and all three
premise relations are well-founded, then for every claim~$x$,
\begin{equation}
  \tauCan_{S''}\bigl((K;L)(x)\bigr)
  \;=\;
  \tauCan_S(x).
\end{equation}
Moreover, the equation factors as the conjunction of the two layer-wise
preservations.
\end{theorem}

\begin{proof}
By lean.
\end{proof}

\begin{remark}
Position maps compose as plain functions between step types---cofunctor
composition---not via a comonadic duplicate. Trust is preserved end-to-end
if and only if it is preserved layer by layer.
\end{remark}

%======================================================================
\section{The shape-only counterexample}
%======================================================================

A consolidation that restates claims but drops checking steps is not a
consolidation in the sense of Definition~\ref{def:consolidation}, and
trust degrades.

\begin{definition}[Ground system]\label{def:ground}
Let $\mathrm{ground}$ be the trust system over a singleton claim set with
a single ground checking step graded~$\mathrm{hi}$ (empty premises).
\end{definition}

\begin{definition}[Shape-only target]\label{def:layer4}
Let $\mathrm{layer4}$ be the trust system over a singleton claim set with
\emph{no} checking steps at all.
\end{definition}

\begin{proposition}\label{prop:ground-layer4}
Both premise relations are well-founded (they are empty). Moreover
\[
  \tauCan_{\mathrm{ground}}(\ast) \;=\; \mathrm{hi},
  \qquad
  \tauCan_{\mathrm{layer4}}(\ast) \;=\; \bott.
\]
\end{proposition}

\begin{proof}
By lean.
\end{proof}

\begin{remark}
No consolidation $\mathrm{ground} \to \mathrm{layer4}$ can exist: there is
no target step for the forward step map to hit. If one did, Theorem
\ref{thm:tau-consolidate} would force
$\tauCan_{\mathrm{layer4}}(\ast) = \mathrm{hi}$, contradicting the
computation above. Shape-only restatement without position pullbacks is
precisely the failure mode that degrades trust from top to bottom.
\end{remark}

%======================================================================
\section{Discussion}
%======================================================================

\paragraph{What has been established.}
Canonical trust is the greatest sound assignment of check grades to
claims, equivalently the max--min grade over the derivation tree, under
well-foundedness of the premise relation. Morphisms transport trust
laxly; consolidations with genuine position maps preserve it exactly and
compose.

\paragraph{Corrections discovered by formalisation.}
Two hypotheses, invisible or underspecified in informal prose, are forced
by the machine-checked development:
\begin{enumerate}
  \item the check-poset must be a chain (for attainment of finite joins
        in the soundness of~$\tau$);
  \item morphisms must send bottom to bottom (for transport on
        underivable claims).
\end{enumerate}

\paragraph{Machine checking.}
Every definition and theorem statement above corresponds to a declaration
in a machine-checked formalisation. The phrase ``by lean'' indicates that
the proof has been verified by a proof assistant. A mathematician reading
this note need only verify that the English statements match the intended
mathematics.

%======================================================================
\section*{Acknowledgements}
%======================================================================

The presentation here is a prose rendering of a formal development
intended for human review of the theorem statements.

\end{document}
